\section{A fixed fourteen-part terminal template proves mixed-source compensation from sixteen onward}
NINE new scoped standard-axiom reports verify the whole contraction,
capacity envelope, block and convolution bounds, terminal product,
support guards and bulk payment below. The actual word/CI/coefficient
identifications and final mixed-head joins are the written assembly,
not an end-to-end Lean composition implementation. There is NO finite
word proof premise. The GENERAL original conjecture and general C remain OPEN.

Keep $\rho=3/5$, $b=35/32$, $q=6125/6144$, $v=175/192$,
$K_0=125/22$, and $s_t=S_t(\rho)$. Put $r_t=s_t/b^{t-1}$.
The frozen norm gives $F(\rho)+K_0\le10b^M$ for EVERY formal word.
For every append $t\ge3$, its stronger exact step is
\[
       F'(\rho)+K_0\le s_t(F(\rho)+K_0),              \tag{1}
\]
because $K_0(s_t-1)\ge1$. Thus an ENTIRE arbitrary tail of capacities
at least three has multiplicative transport of this offset norm.

\subsection{Three new sharp but elementary unbounded scalar bounds}
The ratios $r_t$ decrease for all $t\ge4$, and $r_3\le r_4$.
Indeed $s_{t+1}=s_t+\rho^{t+1}$ and, for $t\ge4$,
\[
 \rho^{t+1}\le\rho^5\le(b-1)s_4\le(b-1)s_t.
\]
Consequently $r_t\le r_{\max(4,L)}$ whenever $t\ge L\ge3$.
For ten arbitrary capacities $u_i\ge i+3$, $i=0,\ldots,9$,
\[
 \prod_{i=0}^9r_{u_i}\le P_*:=r_4^2\prod_{t=5}^{12}r_t.        \tag{2}
\]

A sharper WHOLE block gives
\[
                      m q^m\le128\qquad(m\ge0).       \tag{3}
\]
Use $q^{16}\le20/21$, write $m=16a+r$, $0\le r<16$, and
$(a+1)(20/21)^a\le8$. The latter sequence increases through index 19
and decreases thereafter; alternatively check indices 0..19 exactly
and use the decreasing-ratio induction after 19. This is a fixed
scalar base plus unbounded induction, with NO finite word premise.

For EVERY integer $i\ge0$,
\[
                       (i+1)\rho^i\le6/5.             \tag{4}
\]
Index zero is one, index one is $6/5$, and the consecutive ratio is
at most one after index one. Since $[u][w]$ has coefficient at rank
$i$ at most $i+1$, (4) gives, for every nonnegative polynomial $G$,
\[
              \rho^k[z^k]([u][w]G)\le(6/5)G(\rho).   \tag{5}
\]
This is a WHOLE convolution bound: bound each weighted kernel
coefficient and sum all used nonnegative $G$ terms. It does not
differentiate a coefficient inequality.

The single exact terminal-template constant is
\[
                \frac{41472}{125}P_*\le b^{39}.        \tag{6}
\]
The ratio of the two sides is approximately $0.9764430749$;
only the exact rational inequality, not this decimal, is a proof input.

\subsection{Bulk slope bound from a fixed last fourteen parts}
If an ACTUAL increasing arm has first missing weight at least sixteen,
it has at least fourteen non-one parts, and its LAST fourteen capacities
are at least $2,3,\ldots,15$. This fixed template suffices; no upper
bound on the first missing weight or any capacity is imposed.

Write the last three capacities as $w,u,T$, in nondecreasing order,
and let $G$ be the earlier forward arm, of skeleton degree $M_G$.
The first of the terminal fourteen capacities is absorbed into the
earlier prefix. The following TEN capacities are respectively at least
$3,4,\ldots,12$, so the frozen global norm and (1)--(2) give
\[
                       G(\rho)+K_0\le10b^{M_G}P_*.
                                                               \tag{7}
\]
The three deleted capacities are at least $13,14,15$, so their total
skeleton degree is at least $12+13+14=39$.

Consider any $m$ with $M\le2m$, where $M$ is the whole main-box degree.
The prefix before the LAST TWO parts has skeleton degree at least
$1+1+\cdots+12=79$. Since $T\ge u$, this implies $M\ge2u+4$.
Hence $m-1>u$. The exact forward identities therefore give
\[
 \begin{split}
 \Delta F_Y(m)&\le F_{Gwu}(m-1),\\
 F_{Gwu}&=1+S_u+S_uS_wG,\\
 F_{Gwu}(m-1)&=[z^{m-3}]([u][w]G).
 \end{split}                                                    \tag{8}
\]
Both source terms in the middle line vanish at the used coefficient;
this support guard is essential. Apply (5), (7), and the removed degree:
\[
 b^{39}\frac{\Delta F_Y(m)}{2^m}
     \le10\rho^3(6/5)q^m P_*.
\]
Multiplication by $m$ and (3)--(6), using
$10\rho^3(6/5)\cdot128=41472/125$, prove
\[
                          m\Delta F_Y(m)\le2^m.       \tag{9}
\]
No curvature, C, A1, A2 or source-prefix premise is used for this slope bound.

\subsection{Other-degree branch and all mixed proper heads}
Let $s_Y=\min(\lfloor M/2\rfloor,M-B)$, with $B$ the largest interval
degree. The nonbulk alternative $m>M-B$ is exactly the same as in the
frozen twenty-threshold proof. For first missing at least sixteen,
$M-B\ge1+13\cdot14/2=92$. Replace the old large-read base 100 by 80:
\[
                10(j+1)(175/192)^j\le1\qquad(j\ge80),
\]
whose base is an exact rational check and whose consecutive ratio
decreases thereafter. This covers BOTH a largest last capacity and a
largest initial interval, including the zero-initial-run formal comparison.
Combining with (9) gives EVERY $m>s_Y$ the signed slope bound (9).

The frozen COMPLETE indexed-owner gap and dyadic-prefix/source-sum
theorems then imply every positive mixed recovery has $j+1\le s_Y$.
All properly conditioned CLOSED S heads are paid by the genuine strict
opposite derivative prefix, the one-term absolute budget, the frozen
relative $+1$ surplus and the negative-birth unit argument. Nonpositive
recoveries retain the usual OLD nonnegative head budget.

The accompanying independent reconstruction of the exact constants and
actual support/source identities verifies this assembly. It covers ALL
rising-opposite first missing weights at least sixteen. The remaining
mixed gap is ONLY $K=13,14,15$, with owner weight at most 287 and skeleton degree
at most 143 by the frozen unbounded owner bounds. General higher-capacity
C/both-rising resources, FULL A2 and the ORIGINAL conjecture remain OPEN.
